\usepackage{amsmath}
\usepackage{amsthm}
\usepackage{amssymb}
\usepackage{graphicx}
\usepackage{mathrsfs}
\usepackage{enumitem}
\usepackage{geometry}
\usepackage[colorlinks, linkcolor=black]{hyperref}
\usepackage{stackengine}
\usepackage{yhmath}
\usepackage{extarrows}
% \usepackage{unicode-math}
\usepackage{esint}
\usepackage{multicol} % 导入multicol包
\usepackage{fancyhdr}
\usepackage[dvipsnames, svgnames]{xcolor}
\usepackage{listings}
\usepackage{physics}

\definecolor{mygreen}{rgb}{0,0.6,0}
\definecolor{mygray}{rgb}{0.5,0.5,0.5}
\definecolor{mymauve}{rgb}{0.58,0,0.82}
\definecolor{NavyBlue}{RGB}{0,0,128}
\definecolor{Rhodamine}{RGB}{255,0,255}
\definecolor{PineGreen}{RGB}{0,128,0}

\graphicspath{ {figures/},{../figures/}, {config/}, {../config/} }

\linespread{1.6}


\usepackage{float}    % 提供[H]参数
\usepackage{adjustbox}% 提供垂直对齐控制

% 设置页面布局，包括装订线
% left=左边距, right=右边距, top=上边距, bottom=下边距
% bindingoffset=装订线宽度，通常在1cm到1.5cm之间
\geometry{
    a4paper,
    inner=3cm,     % 内侧边距（含装订线）
    outer=2cm,     % 外侧边距
    top=2.5cm,
    bottom=2cm,
    bindingoffset=1cm  % 从内侧边距中划分出装订线
}
\setenumerate[1]{itemsep=5pt,partopsep=0pt,parsep=\parskip,topsep=5pt}
\setitemize[1]{itemsep=5pt,partopsep=0pt,parsep=\parskip,topsep=5pt}
\setdescription{itemsep=5pt,partopsep=0pt,parsep=\parskip,topsep=5pt}

\lstset{
    language=Mathematica,
    basicstyle=\tt,
    breaklines=true,
    keywordstyle=\bfseries\color{NavyBlue}, 
    emphstyle=\bfseries\color{Rhodamine},
    commentstyle=\itshape\color{black!50!white}, 
    stringstyle=\bfseries\color{PineGreen!90!black},
    columns=flexible,
    numbers=left,
    numberstyle=\footnotesize,
    frame=tb,
    breakatwhitespace=false,
} 

\lstset{
    language=TeX, % 设置语言为 TeX
    basicstyle=\ttfamily, % 使用等宽字体
    breaklines=true, % 自动换行
    keywordstyle=\bfseries\color{NavyBlue}, % 关键字样式
    emphstyle=\bfseries\color{Rhodamine}, % 强调样式
    commentstyle=\itshape\color{black!50!white}, % 注释样式
    stringstyle=\bfseries\color{PineGreen!90!black}, % 字符串样式
    columns=flexible, % 列的灵活性
    numbers=left, % 行号在左侧
    numberstyle=\footnotesize, % 行号字体大小
    frame=tb, % 顶部和底部边框
    breakatwhitespace=false % 不在空白处断行
}
% 定义注意环境
\newenvironment{note}{%
  \par\noindent% 开始新段落，取消首行缩进
  \textbf{注意：}% 注意前缀
  \begin{list}{%
    \arabic{enumi}.% 列表编号格式：1.、2.
  }{%
    \usecounter{enumi}% 使用默认计数器
    \setlength{\labelwidth}{1.2em}% 编号宽度
    \setlength{\labelsep}{0.3em}% 编号与文字间距
    \setlength{\leftmargin}{\dimexpr\labelwidth + \labelsep\relax}% 整体缩进
    \setlength{\itemindent}{0pt}% 取消项缩进
    \setlength{\topsep}{0pt}% 取消顶部间距
    \setlength{\partopsep}{0pt}% 取消顶部额外间距
    \setlength{\parsep}{0pt}% 取消段落间距
    \setlength{\itemsep}{0.5em}% 列表项之间的间距
  }%
  \let\olditem\item% 保存原始\item命令
  \renewcommand{\item}{% 重定义\item命令
    \ifnum\value{enumi}=0\relax% 判断是否为第一个列表项
      \olditem% 第一个项正常显示（与"注意："同行）
    \else%
      \olditem% 后续项保持正常缩进
    \fi%
  }%
}{%
  \end{list}% 结束列表环境
  \par% 结束段落
}
% \begin{lstlisting}[language=TeX] ... \end{lstlisting}

\renewcommand{\eqref}[1]{%
  \hyperref[#1]{式\ref*{#1}}% 
  % \ref* 避免重复生成跳转链接（仅 \hyperref 负责跳转）
}